\chapter{资产管理与机构投资}
\label{chap:Asset_Management_and_Institutional_Investment}

\begin{openingcase}{管理人为谁的利益服务？}

一个养老基金代表数千名受益人进行退休储蓄的投资。这些受益者最终承担投资带来的收益和损失，但他们不会单独选择证券或自行交易投资组合。相反，该基金将大部分责任委托给专业的资产管理者来负责。

有一位管理者发现了一种投资机会，这种投资方式需要大幅偏离基准水平，且可能需要数年时间才能获得回报。如果这项投资成功，那么收益主要将归那些享受养老金福利的人所有。然而，如果这项投资在短期内表现不佳，那么养老基金可能会削减或终止对该管理者的委托责任，这样一来，该管理者的费用收入也会相应减少。因此，这位管理者可能会犹豫是否要采取符合受益人长期利益的决策，因为这样的决策会带来短期内的风险。

\end{openingcase}


\begin{tcolorbox}[
  enhanced,
  colback=FinanceBlue!6,
  colframe=FinanceBlue,
  boxrule=0.6pt,
  sharp corners,
  left=8pt,
  right=8pt,
  top=7pt,
  bottom=7pt,
  before skip=1em,
  after skip=1.6em,
  title=\textbf{核心问题},
  coltitle=white,
  colbacktitle=FinanceBlue,
  fonttitle=\bfseries,
  attach boxed title to top left={xshift=6pt,yshift=-2mm},
  boxed title style={
    boxrule=0pt,
    sharp corners,
    left=5pt,
    right=5pt,
    top=2pt,
    bottom=2pt
  }
]
\begin{enumerate}[
  leftmargin=*,
  label=\textcolor{FinanceBlue}{\textbf{\arabic*.}},
  itemsep=0.55em,
  topsep=0.2em
]


\item 为什么投资者会委托进行投资组合决策？而指令要求、基准指标、费用以及信托责任又如何影响着投资者与资产管理者之间的关系呢？

\item 不同的资产管理工具和机构投资者在流动性、控制权、投资期限以及风险分配方面有何差异？

\item 资产管理者如何生成信息？他们如何影响市场以及企业治理？同时，他们如何应对被动投资和技术驱动的投资呢？


\end{enumerate}
\end{tcolorbox}

\section{委托资产管理的经济逻辑}

资产管理将资本所有权与许多投资决策分离开来。投资者承担投资组合的盈亏，而专业管理者则有权在规定的职责范围内选择、交易、监控证券，有时甚至有权投票决定证券的处置方式。

\subsection{直接投资与委托投资}

个人投资者可以直接持有和管理金融资产，或者将投资组合的决策委托给专业的基金经理。直接投资让投资者能够掌控证券的选择、交易以及投资组合的监控，但这需要时间、信息以及投资方面的专业知识。而在委托投资的情况下，投资者仍然对投资组合承担经济责任，而基金经理则有权在规定的授权范围内做出具体的投资决策。具体的法律安排取决于资产是通过集合平台、单独账户还是其他形式来持有的。

当投资组合管理需要专业知识、规模或基础设施时，委托给专业人士来管理可以创造价值。对于个人投资者来说，这些资源往往难以获得。管理者可以将研究、交易系统、风险测量、报告和管理等成本分摊到多个客户身上。通过汇集投资，还可以以比单个投资者更低的交易和监控成本来实现多元化投资。

这一结构不同于吸收存款的中介。商业银行主要通过在自身资产负债表上发行存款和其他负债，融资其资产组合。资产管理人则通常管理客户资产或集合基金，而非主要用自身负债融资这些资产。因此，组合损益通常由投资者承担，但具体仍取决于载体法律结构，以及管理人提供的契约承诺或资本。

一个简单的委托投资关系可以表示为：

\begin{equation}
W_T=W_0(1+r_{gross})-F,
\label{eq:ch14-investor-terminal-wealth}
\end{equation}
其中，$W_0$ 是投资者的初始资本，$r_{gross}$ 是期间投资组合的毛回报率，$F$ 则代表总费用和其他费用，以货币金额表示。与直接投资相比，委托投资改变了做出投资组合决策的人选，并引入了管理成本，而投资者通常仍然需要承担投资组合的业绩风险。

\subsection{资产管理契约、费用与激励}

委托资产管理是由定义决策权、投资限制和补偿条款的合同来管理的。投资指令可能会规定适用的资产、基准指标、风险限额、杠杆率、投资期限以及报告要求。在集合投资工具中，基金条款还可能规定流动性及提款的相关事宜。所有这些规定共同构成了管理层的决策空间以及评估绩效的依据。

基准提供了一个参考投资组合或回报指标，可以通过该指标来评估投资组合的表现。如果$r_P$表示投资组合的回报，$r_B$则表示基准回报，那么主动回报就是……

\begin{equation}
r_{active}=r_P-r_B.
\label{eq:ch14-active-return}
\end{equation}

基准测试能够更清晰地显示相对性能和投资组合的偏差，但这也可能影响投资决策。如果管理者始终以基准作为衡量标准，可能会限制那些可能导致显著相对表现不佳的较大偏差。

补偿通常包括基于资产的管理费用。如果年度管理费率为$f$，而$A$则代表相关的费用基数，那么管理费用可以表示为：

\begin{equation}
F_M=fA.
\label{eq:ch14-management-fee}
\end{equation}

一些安排还包括基于绩效的补偿机制。简化的绩效奖励制度如下：

\begin{equation}
F_P=\alpha\max\{r_P-h,0\}A,
\label{eq:ch14-performance-fee}
\end{equation}
其中，$\alpha$表示绩效费用率，$h$表示门槛值或参考收益水平，$A$则表示费用的基础金额。更详细的合同条款可能会包括最高限额、损失承担条款，或者基准调整相关的内容。

费用设计影响管理人激励。按资产规模收费，使收入与管理资产规模相联系，提高吸引和留住投资者资金的价值。业绩报酬使报酬对投资结果更敏感，但在某些条件下，不对称参与收益与损失也可能鼓励更高风险承担。最终激励取决于整个契约，以及管理人是否与客户共同投入资本。

\begin{example}{管理费与投资者净回报}{management-fee-net-return}

该基金年初拥有1亿元人民币的资金，其投资组合的毛回报率为8\%。假设管理费的支出为初始资产的1\%，且没有其他费用或现金流。

管理费为

\[
F_M=0.01\times100=1
\]

投资组合的收益使基金的资产在扣除费用前增加到1.08亿元人民币。扣除管理费后，净资产为1.07亿元人民币，对应的简单净回报率为7\%。

总回报与净回报之间的每百分点差异，实际上反映了按照既定惯例所计算出的管理费用。

\end{example}

\subsection[委托代理问题与受托责任]{委托代理问题与受托责任}

委托资产管理会形成一种本金与代理人的关系。投资者将投资权限委托给管理者，但管理者的目标并不总是完全符合投资者的利益。补偿、职业考量、资产积累以及商业关系都可能影响管理者行使权力的方式。

机构之间的利益冲突可以有多种表现形式。当补偿机制不对称时，管理者可能会承担额外的风险；为了降低职业发展的风险，他们可能会坚持遵循某个基准标准；或者出于某种考虑，他们可能会进行超出投资合理范围的交易行为。这些行为并非委托授权过程中固有的现象；其重要性取决于合同设计、组织结构以及监管环境等因素。

受托责任约束委托权力的使用。其法律内容因司法辖区和机构形式而异，但基本经济原则是：管理人应充分考虑客户利益，行使受托权力，并适当管理利益冲突。披露、内部控制、独立监督和对利益冲突交易的限制，均可支持这一目标。

投资者还可以通过监控和退出来约束基金经理的行为。机构投资者可以协商投资指令，评估基金经理的表现，并更换基金经理；而散户投资者则更多地依赖标准化的信息披露、基金治理机制以及赎回或转换基金的能力。当未来的投资指令和资产流动依赖于投资者的信心时，声誉因素就成为了一个额外的制约因素。

\begin{tcolorbox}[
  enhanced,
  breakable,
  colback=white,
  colframe=FinanceBlue,
  boxrule=0.7pt,
  sharp corners,
  borderline west={2.2pt}{0pt}{FinanceBlue},
  title=\textbf{财务展望：诚信责任并不能消除代理问题},
  coltitle=white,
  colbacktitle=FinanceBlue,
  fonttitle=\bfseries,
  before skip=1.2em,
  after skip=1.2em
]

受托责任约束管理人如何使用委托权力，却不会使投资者与管理人的激励完全一致。投资结果不确定，管理行为未必完全可观察，利益冲突也可能持续存在。因此，即使存在受托义务，契约、监督、治理和声誉仍然重要。

\end{tcolorbox}

\section{集合投资载体}

一种常见的组织委托投资的方式是通过集体投资机构。这种机构不是单独管理每个投资者的投资组合，而是汇集许多投资者的资金，并将这些资源共同用于特定的目标。

\subsection{开放式共同基金}

开放式共同基金发行投资者可向基金申购或赎回的份额。申购创造新份额，赎回注销现有份额，因此已发行在外份额数量随投资者资金流动而变化。

每股资产净值是多少？

\begin{equation}
\textit{NAV}_t=
\frac{A_t-L_t}{N_t},
\label{eq:ch14-nav}
\end{equation}

其中，$A_t$~为基金资产市值，$L_t$~为基金负债，$N_t$~为已发行在外份额数量。单位净值将基金组合净价值换算为每份价值。

投资者流动可能会影响投资组合的管理。认购资金可以用于投资更多的资产，而赎回则要求基金向离开的投资者提供现金。基金可以满足因持有现金而产生的赎回需求，同时还可以用于弥补认购资金、允许借款或出售资产。当赎回金额相对于可用流动性而言较大时，即使基础投资的状况没有变化，基金经理也可能需要出售资产。

开放式资金池可以提供多样性和交易管理的便利性，但同时也使得投资者的流动性与投资组合的流动性之间产生了联系。因此，满足赎回需求的成本部分取决于基础资产的可销售性。

\begin{example}{按单位净值申购共同基金}{mutual-fund-subscription-nav}

该共同基金的资产为人民币约5亿元，负债为人民币约500万元，存续股份数量为4500万股。其每股净资产值为

\[
\textit{NAV}=\frac{500-5}{45}=11.00.
\]

一位投资者以110万元人民币的出资参与了该项目的投资。扣除相关费用后，投资者实际获得的收益为……

\[
\frac{1.1\text{ 百万}}{11.00}=100{,}000
\]

新发行的基金份额使基金资产增加至5.011亿元人民币，存续份额达到4510万份。按照既定的假设，每股净资产价值仍保持在11元人民币左右。

\end{example}

\subsection[交易所交易基金与创设赎回机制]{交易型基金与创建——\\赎回机制}

交易所交易基金（ETF）持有基础资产组合，其份额在二级市场交易。多数投资者买卖~ETF~份额时，不直接与基金交易。但授权参与人可通过与基金交易，大批量创设或赎回份额。

创设时，授权参与人交付指定的一篮子证券、现金或两者组合，获得新发行的~ETF~份额；赎回时，将~ETF~份额交回基金，换取相应资产篮子或现金结算。这一机制使~ETF~份额供给能够根据其自身及基础组合的市场条件调整。

让 $P_t^{ETF}$ 表示 ETF 股票的市场价格，而 $\textit{NAV}_t$ 则表示其每股的资产净值。一个简单的溢价或折价衡量方法就是……

\begin{equation}
\pi_t=
\frac{P_t^{ETF}-\textit{NAV}_t}{\textit{NAV}_t}.
\label{eq:ch14-etf-premium-discount}
\end{equation}

ETF~以足够大的溢价交易时，授权参与人可创设新份额并在二级市场出售，促使溢价缩小；存在足够大的折价时，赎回可发挥相反作用。这一机制的有效性取决于交易成本、融资条件、税收、交易时间和基础资产流动性。

\begin{example}{ETF溢价交易时的份额创设}{etf-creation-premium}

假设一只ETF的市值为每股10.00人民币，但在二级市场上交易价格为10.10人民币。忽略交易成本和其他摩擦因素，授权参与者可以以相当于每股10.00人民币的价格购买基础资产篮子，将篮子交付给基金，然后获得新创建的ETF股份。

如果这些股票以10.10人民币的价格出售，那么价格差异就会创造出套利机会。这种机会会增加ETF股票的供应，从而有助于降低溢价率。

实际上，只有当价格差异足够大以覆盖交易成本、融资成本、税收以及其他摩擦时，创建才变得有吸引力。

\end{example}

不同于普通开放式共同基金，ETF~不要求每次投资者买卖份额时基金自身都进行交易。因此，二级市场可以吸收大量投资者交易，而不直接改变基础组合。但创设赎回机制仍为交易压力在~ETF~份额与基础资产之间传递提供渠道。

\subsection{基金流动性、单位净值与赎回风险}

请求权投资组合的市场流动性与其底层资产的流动性虽然相关，但又是不同的。即使某些投资组合资产的交易频率较低，基金仍可能提供频繁的申购和赎回操作。当投资者的流动性与可用现金和市场流动性相比较为有限时，这种安排更容易维持。

大量赎回可能要求出售资产。组合资产相对缺乏流动性时，出售可能产生交易成本、价格让步或市场冲击。如果赎回投资者获得的价格没有充分反映这些成本，部分成本可能由留在基金中的投资者承担。在这种条件下，对进一步资金流出的预期可能形成提前赎回激励。

基金的设计方式会影响这些成本的分配。现金缓冲和流动性储备可以减少对立即出售资产的需求，而赎回费、变动定价、通知期限以及其他流动性管理工具则可以将更多成本转移给那些产生这些成本的投资者。这些工具的运用取决于基金的结构、资产的流动性状况以及相关的规章制度。

流动性不足的资产也会使证券的净值计算变得复杂。当市场价格无法获取或获取频率较低时，估值可能会更多地依赖于模型、交易商的报价、可比证券或其他估算方法。虽然证券的净值仍然是衡量投资者请求权的一个有用指标，但当可观测的市场价格提供的信息较少时，估值的不确定性可能会增加。

\begin{figure}[H]
\centering
\begin{tikzpicture}[
  >=Latex,
  box/.style={
    draw=FinanceBlue,
    rounded corners=2pt,
    align=center,
    text width=3.0cm,
    minimum height=0.9cm,
    inner sep=5pt
  },
  flow/.style={-Latex, line width=0.8pt, FinanceBlue}
]

\node[box] (redeem) at (0,0) {投资者赎回资金};
\node[box] (cash) at (4.3,0) {需要现金。};
\node[box] (sales) at (8.6,0) {投资组合资产销售};

\node[box] (costs) at (8.6,-2.0) {交易成本与市场影响};
\node[box] (remain) at (4.3,-2.0) {由其余投资者部分承担的成本};
\node[box] (incentive) at (0,-2.0) {可能的激励措施，以提前实现兑换目标。};

\draw[flow] (redeem) -- (cash);
\draw[flow] (cash) -- (sales);
\draw[flow] (sales) -- (costs);
\draw[flow] (costs) -- (remain);
\draw[flow] (remain) -- (incentive);

\end{tikzpicture}
\caption{开放式基金的赎回压力与流动性成本}
\label{fig:ch14-redemption-liquidity}
\end{figure}

\section{机构投资者与另类资产管理}

除了集体投资工具之外，委托投资还通过各种具有不同负债结构、投资期限、流动性需求以及治理安排的机构来组织实施。这些特征影响着它们能够持有的资产、能够承担的风险，以及它们与金融市场和企业发行人的互动方式。

\subsection{养老基金与负债驱动型投资}

养老基金投资资产的目的是为了履行未来的福利义务。因此，它们的投资策略取决于这些支付的时机、通货膨胀的影响程度以及支付金额的不确定性，同时也需要考虑预期的资产回报情况。

设预期的养老金支付在日期$t=1,\ldots,T$时为$B_t$。负债的现值简化表示为

\begin{equation}
L_0=\sum_{t=1}^{T}\frac{B_t}{(1+y_t)^t},
\label{eq:ch14-pension-liability-pv}
\end{equation}
其中，$y_t$ 表示应用于第 $t$ 天支付的贴现率。衡量资金状况的一个简单方法是……

\begin{equation}
FR=\frac{A_0}{L_0},
\label{eq:ch14-pension-funding-ratio}
\end{equation}
其中，$A_0$ 表示养老金资产的市场价值。

负债驱动型投资（LDI）相对于养老金负债的价值和敏感性管理资产。例如，长期债券的久期若更接近负债久期，可减少利率变化对资金充足状况的影响。基金也可能持有其他资产，以寻求更高预期回报。具体配置取决于资金充足状况、给付设计、缴费政策、监管约束，以及对资金充足程度波动的容忍度。

长期负债还可以使养老基金在短期福利支付和抵押品需求仍然可控的情况下，拥有持有流动性较低的资产更大的能力。虽然较长的投资期限并不能消除流动性风险；但当有足够的流动性来满足短期债务需求时，它可以扩大可行的投资范围。

\begin{tcolorbox}[
  enhanced,
  breakable,
  colback=white,
  colframe=FinanceBlue,
  boxrule=0.7pt,
  sharp corners,
  borderline west={2.2pt}{0pt}{FinanceBlue},
  title=\textbf{财务前景：LDI专注于资金状况，而不仅仅是资产状况。},
  coltitle=white,
  colbacktitle=FinanceBlue,
  fonttitle=\bfseries,
  before skip=1.2em,
  after skip=1.2em
]

负债驱动型投资相对于养老金负债评估资产风险。资产组合单独看可能表现良好，但若负债价值上升更多，资金充足状况仍会恶化。因此，关键在于资产与负债价值如何共同变化。

\end{tcolorbox}

\subsection{主权财富基金、捐赠基金与长期资本}

主权财富基金、捐赠基金、基金会以及类似的机构通常会在较长的时期内进行投资，但它们的资本来源和支出义务却存在显著差异。主权财富基金可能通过商品收入、财政盈余、外汇储备或代际储蓄来融资，而捐赠基金则通常用于支持机构的经常性支出，并力求在一段时间内保持购买力的稳定。

这些差异会影响流动性需求和资产配置。较长的投资期限可以增加持有流动性较低的资产的能力，比如私募股权、基础设施和房地产等，前提是短期现金需求仍然可控。因此，仅仅拥有较长的投资期限是不够的；该机构还必须具备维持长期头寸的流动性和管理能力。



\begin{tcolorbox}[
  enhanced,
  breakable,
  colback=white,
  colframe=FinanceBlue,
  boxrule=0.7pt,
  sharp corners,
  borderline west={2.2pt}{0pt}{FinanceBlue},
  title=\textbf{财务前景：长远的规划并不能消除流动性方面的限制问题},
  coltitle=white,
  colbacktitle=FinanceBlue,
  fonttitle=\bfseries,
  before skip=1.2em,
  after skip=1.2em
]

较长的投资期限可以增加机构持有流动性较低资产的能力，但仅凭期限是不够的。支出需求、可能的提款、抵押品要求以及治理限制仍然可能需要现金支持。因此，长期投资的能力取决于投资期限和流动性能力。

\end{tcolorbox}

支出规则和机构目标进一步影响着投资组合的选择。拥有相对稳定年度收益的捐赠基金可能会更重视长期实际回报，而具有稳定目标的主权基金在不利经济条件下 withdrawals 可能会增加流动性资产的持有量。

长期资本投资还可以降低在市场出现短期压力时需要进行强制出售的需要。拥有足够流动性的机构能够在价格出现短暂下跌时继续持有头寸，而这样做仍然符合它们的职责要求以及风险承受能力。

\subsection{私募股权与风险投资}

私募股权与风险投资通常通过封闭式基金组织投资，有限合伙人承诺投入资本，普通合伙人作出投资决策。投资者一般无法随时赎回权益，使基金有更大能力长期持有非流动性投资。

私人市场投资的一个核心特点是集中所有权和实时监控。收购基金可能会获得实质性的控制权，从而能够影响公司的融资、治理、管理和运营决策。风险资本投资者通常会分阶段为年轻企业提供资金支持，并可能通过参与董事会和管理层来进一步干预公司事务。在信息有限的情况下，这种合作方式尤为重要，因为企业的价值很大程度上取决于其未来的发展情况。

管理人报酬通常由管理费与挂钩投资收益的业绩分成构成，后者常称附带收益。这使普通合伙人直接参与投资结果，但仍可能存在激励冲突。基金条款、分配规则和管理人投入的资本，都会影响普通合伙人与有限合伙人如何分享收益、损失与激励。

当投资组合中的公司并非持续被交易时，估值就变得更加困难。基金的价值可能依赖于模型、可比交易或其他在可观察市场交易中得出的估值方法。因此，报告的估值可能会比公开市场的价格调整得更慢，并且存在更大的估计不确定性。

\subsection{对冲基金与另类投资策略}

对冲基金通常遵循相对灵活的投资策略，可能会使用卖空、衍生品、杠杆、相对价值头寸、事件驱动策略或集中性持仓。这些工具使得基金经理能够表达正面和负面的观点，对冲选定的风险，并构建在传统长期持有投资组合中难以实现的头寸。

这些策略可能会引入资金风险、利润率风险、流动性风险以及模型风险，此外还有基础头寸的市场风险。因此，这些风险不仅取决于所持有的资产，还取决于如何为这些头寸提供融资、进行组合操作以及采取何种抵押措施。

投资者流动性条款通常旨在反映策略的流动性状况。锁定期、限制措施等规定可能会减少因提款而需要进行强制资产出售的情况，但同时也会限制投资者快速退出市场的能力。

权益可以放大资金权益的变化。如果一个基金的资产$A$部分由权益$E$来融资，那么一个简单的资产负债表杠杆指标就是$A/E$。

这一比率并未涵盖所有形式的经济杠杆效应。衍生品、空头头寸以及其他风险敞口可能会产生一些在资产负债表资产中未能完全反映的敏感性问题。因此，风险分析还需要考虑投资组合的风险敞口、抵押品要求以及融资条件。

\section[作为信息与治理中介的资产管理人]{作为信息与治理中介的资产管理人}

一旦投资决策被委托出去，资产管理者就可以影响不仅仅是投资组合的分配。通过研究、交易、投票和参与等方式，他们可以影响市场价格所反映的信息，以及他们所持有证券的公司治理。

\subsection{信息生产、研究与证券选择}

活跃的资产管理团队会收集和使用信息，以确定那些预期回报相对于价格和风险来说具有吸引力的证券或投资组合。研究可能会参考财务报表、行业和宏观经济状况、公司信息、定量模型以及其他允许的数据来源。其价值取决于是否能够在价格完全确定之前获取相关信息并将其纳入投资决策中。

信息的生产与实施成本很高。研究需要分析师、数据和技术支持，而投资观点则必须通过交易转化为投资组合中的具体仓位。随着其他市场参与者获取并利用类似的信息，这种信息优势的价值可能会下降。

如果$r_{gross}$表示在管理费和实施成本之前的投资组合回报，$f$和$c$分别表示相应的成本率，那么简化的净回报就是

\begin{equation}
r_{net}=r_{gross}-f-c.
\label{eq:ch14-net-portfolio-return}
\end{equation}

只有当信息、证券选择以及执行方面的贡献足以抵消这些成本时，主动管理才能为投资者带来价值。

规模可以以不同的方式影响这一过程。研究、数据和技术可能具有规模经济效应，而大型投资组合在资源有限的情况下进行部署时可能会更加困难，或者在进行交易时可能会影响价格。因此，所管理的额外资产的价值取决于其运营市场的策略以及市场的流动性。

\subsection{机构持股与公司治理}

机构投资者可以通过积极的监督以及调整自己的投资布局来影响公司治理结构。这些渠道的有效性取决于企业的持股规模、投票权分配、股权集中度，以及投资管理的相关规定。

管理指的是代表客户或受益人积极监督被投资公司的行为。根据客户的投资需求，资产管理者可以行使与客户投资组合中的股份相关的投票权，包括对董事选举、高管薪酬、合并以及其他股东决议的投票权。他们还可以与董事会和管理层就公司战略、业绩和治理问题进行沟通。通过协调这些活动，资产管理者可以在股东监督方面实现规模经济，并有助于解决与分散式股份所有权相关的集体行动问题。

机构持股也可能引发治理问题。大型资产管理公司在许多公司中拥有巨大的投票影响力，而与被投资公司的业务关系或客户利益之间的冲突则可能引发利益冲突。同时，对竞争企业的共同持股也引起了关注，因为机构持股的重叠可能会影响企业之间的竞争动机。这些问题的严重性仍有争议，不同市场以及不同的所有权结构下，实际情况也有所不同。

\begin{tcolorbox}[
  enhanced,
  breakable,
  colback=white,
  colframe=FinanceBlue,
  boxrule=0.7pt,
  sharp corners,
  borderline west={2.2pt}{0pt}{FinanceBlue},
  title=\textbf{财务视角：被动投资并不排除主动管理},
  coltitle=white,
  colbacktitle=FinanceBlue,
  fonttitle=\bfseries,
  before skip=1.2em,
  after skip=1.2em
]

指数基金遵循基于基准的投资策略，限制证券选择及组合权重的裁量，但并不妨碍管理人行使投票权、与被投资公司沟通或监督公司治理。仅因治理担忧卖出指数成分证券，可能产生跟踪误差，因此投票和沟通可成为被动机构投资者发挥影响的重要渠道。

\end{tcolorbox}

\subsection{机构交易与市场流动性}

机构投资者交易证券是为了实施投资组合决策、重新平衡投资组合、适应投资者流动以及管理风险。他们的订单在其他交易的一方时可以提供流动性，而在寻求即时执行时则可能需要流动性。

大型订单需要在执行速度和市场影响之间做出权衡。较快的执行速度可以降低价格波动带来的风险，但此时订单可能无法完全执行，因此需要更高的价格让步。相反，较慢的执行速度虽然可以减少瞬间的市场影响，但却会增加在执行过程中面临价格变动的风险。

实施不足（IS）衡量的是与投资决策时现行价格相对的执行成本。对于完全执行的购买行为，令$P_0$表示决策价格，$\bar{P}$表示平均执行价格，$Q$表示购买的数量。忽略佣金和其他交易成本。

\begin{equation}
IS=Q(\bar{P}-P_0).
\label{eq:ch14-implementation-shortfall}
\end{equation}

正值表示平均执行价格高于决策价格，这意味着相对于该基准值而言，存在执行成本。而在出售的情况下，符号则相反。

机构交易可能会在以下情况下影响市场流动性：当大额或相关性较高的订单超过当前价格下的可成交深度时。由此产生的价格变动取决于订单规模、执行策略、市场深度以及其他参与者是否愿意采取相反的行动。

\section{资产管理的转型}

新的投资结构和技术正在改变资产管理的运作方式。它们可以降低投资组合构建和执行的成本，同时将一些决策从自由裁量权转移到规则和模型之上。

\subsection{被动投资、指数化与基于规则的配置}

被动投资的目标是复制特定基准指标的业绩表现，而不是通过自主选择证券来创造超额回报。指数基金和被动ETF是常见的实现方式，不过类似的策略也可以用于独立管理的账户和机构投资组合中。

基于规则的配置，用预先设定的标准决定组合风险暴露或资产权重。指数化属于规则投资，但并非所有规则策略都是被动策略。一些策略会根据特定市场条件、投资信号或配置规则，系统地调整权重或风险暴露。

一个简单的例子是一个固定权重的投资组合，其中60\%的资金用于股票，40\%用于债券。随着市场价格的变化，实际权重可能会偏离这些目标。一个预先确定的重新平衡规则要求管理者定期或在偏差超过指定阈值时恢复原来的权重。其他基于规则的策略可以根据预定义的波动性目标、估值指标或动量信号来调整投资组合的持仓比例。这些方法减少了对主观决策的依赖，不过在选择指标、阈值和重新平衡规则方面仍然需要投资方面的判断。

被动策略可以在降低研究和投资组合管理成本的同时，为市场或因素提供透明的投资方向。被动投资组合的回报与其基准回报之间的差距被称为跟踪差异。该差异是根据方程~\eqref{eq:ch14-active-return}中的主动回报定义来衡量的。负的跟踪差异表明相对于基准的业绩表现不佳。这可能是由于管理费用、税收、交易成本、现金持有量、抽样以及重新平衡过程中的摩擦所导致的，而证券借贷收入则可能抵消部分这些影响。

被动投资限制了持续证券选择的自由裁量权，但并未消除投资选择的权利。证券提供商通过关于合格证券、纳入标准以及权重分配方法的决策来确定基准构成。资产管理公司则负责选择如何复制这一基准、执行交易、管理现金持有量以及控制跟踪偏差。因此，投资自由裁量权实际上受到基准方法的限制，并且集中在投资组合的运作上。

\subsection{智能投顾与算法投资组合管理}

Robo-advisory利用数字系统将投资者信息转化为投资组合建议，并自动化部分实施和再平衡过程。关于投资期限、风险承受能力、收入需求以及其他限制条件的信息会被转化为资产配置以及一系列投资工具。

自动化可以在这些任务标准化的情况下，降低构建和维护多样化投资组合的成本。它还可以使推荐意见更好地反映更广泛的投资者特征。最终建议的质量取决于收集到的信息、用于解释这些信息的方法，以及系统中内置的约束条件。

算法投资组合管理将这些方法的应用范围扩展到了零售建议之外。资产管理者可以在研究、投资组合构建和交易过程中使用优化算法、执行算法、机器学习技术以及庞大的数据集。这些工具可以提升信息处理能力，并自动化部分投资管理工作。不过，投资目标、风险限制以及法律责任仍然制约着这些工具的广泛应用。

\subsection{模型风险、自动化与人工监督}

基于模型的资产管理将数据、假设和决策规则转化为投资组合的选择。当模型中包含的信息或关系无法充分反映投资问题时，或者当模型的目标与投资者的实际目标不一致时，就会产生模型风险。

自动化可以提升基于模型的决策的实施速度和规模。但它也可能放大错误或常见的反应。如果许多投资组合对相同的信号做出相同的反应，那么在市场流动性有限的情况下，交易压力可能会集中。这种影响取决于投资组合的规模、市场深度以及所使用的策略的多样性。

人工智能能够扩展信息处理的能力，并自动化研究、预测、分类以及通信等过程的部分工作。这些应用或许能够提高效率，但其可靠性仍取决于数据的质量、模型的设计以及输出结果的验证。

人类监督仍然非常重要，因为投资目标和可接受的风险风险是无法通过自动化系统来确定的。监督工作包括明确投资目标、设定风险界限、审查模型假设和输出结果、识别异常情况，以及决定在何种情况下需要人工干预以纠正自动化的规则设置错误。

\begin{table}[H]

\centering

\caption{资产管理自动化：潜在收益与仍存约束}

\label{tab:ch14-automation-asset-management}

\begin{tabularx}{0.94\textwidth}{@{}p{0.30\textwidth} X X@{}}

\toprule

\textbf{函数} & \textbf{自动化带来的潜在收益} & \textbf{剩余的约束条件} \\

\midrule

客户端评估  
& 处理成本更低，一致性更高  
& 客户端信息不完整或不准确  
\\

研究
& 大规模数据集的处理效率提升
& 数据质量与模型可靠性 \\

投资组合构建
& 分配和再平衡规则的一致应用
& 目标明确与估算误差 \\

交易方式  
& 更快且更系统的执行流程  
& 市场影响、流动性以及不断变化的市场状况  
\\

\bottomrule

\end{tabularx}

\end{table}

%========================================================
% Summary
%========================================================

\TailModule{summary}{本章小结}

资产管理将投资组合决策从投资者手中委托给专业管理者。这种委托可以发挥规模经济和信息生产、交易、分散投资和管理的专业优势，但同时也带来了代理问题。投资指令、基准指标、费用、受托责任以及治理安排有助于使管理者的决策与投资者的利益保持一致。通常情况下，投资者仍然需要承担投资组合的盈亏风险。

集体投资工具在发行、交易和赎回请求权方面有所不同。传统的开放式共同基金以资产净值为基础来发行和赎回股份，而ETF则会在交易所进行交易，并将二级市场的交易与创建—赎回机制相结合。因此，基金份额的流动性不仅取决于交易安排，还取决于基础资产的流动性以及交易和赎回成本的分配。

机构投资者在负债、投资期限、流动性需求以及治理结构方面存在差异。养老基金会根据未来的福利义务来管理资产；主权财富基金和捐赠基金通常具有更长的投资期限；私募股权和风险资本基金则结合了承诺的资金与积极的投资决策；而套期保值基金则采用更为灵活的策略，包括使用杠杆和做空交易。

资产管理者也会产生信息，行使所有权，并参与证券市场。被动式、基于规则的和自动化的策略改变了这些功能的执行方式，而不是完全消除它们。在投资决策变得越来越系统化且依赖技术的情况下，模型设计、市场流动性、激励机制以及人工监管仍然非常重要。

%========================================================
% Key Terms
%========================================================

\TailModule{keyterms}{关键术语}

资产管理（Asset management）；委托投资（delegated investment）；投资授权（investment mandate）；委托代理问题（principal-agent problem）；受托责任（fiduciary duty）；管理费（management fee）；基准（benchmark）；主动回报（active return）；尽责管理（stewardship）；共同基金（mutual fund）；单位净值（net asset value）；交易所交易基金（exchange-traded fund）；赎回风险（redemption risk）；负债驱动型投资（liability-driven investing）；资金充足率（funding ratio）；私募股权（private equity）；风险投资（venture capital）；对冲基金（hedge fund）；机构持股（institutional ownership）；被动投资（passive investing）；指数基金（index fund）；基于规则的配置（rules-based allocation）；跟踪偏离（tracking difference）；执行缺口（implementation shortfall）；智能投顾（robo-advisory）；模型风险（model risk）。

%========================================================
% Financial AI Lab
%========================================================

\TailModule{aiprompt}{金融~AI~实验室}

\textbf{实验室：审核由人工智能生成的基金推荐方案}

请一个AI系统来比较两款投资于同一大类资产类别的投资基金，并推荐其中一款。要求该系统解释该基金的架构、基准指标、费用、流动性、投资组合构成，以及相关的历史跟踪差异和主要投资风险。

将建议分为四个阶段进行审核。

\begin{enumerate}

\item \textbf{请检查资金结构。}

请判断每辆汽车是否属于开放式共同基金、ETF、封闭式基金、私人基金，还是其他类型的投资结构。同时，请确认投资者是如何参与和退出投资的，以及交易是否是在市场价格进行的，还是基于净资产值进行的。

\item \textbf{请查看相关基准测试和成本信息。}

请使用官方的资金披露文件来确认所声称的基准值、管理费用以及其他相关费用。如果使用一致的周期和回报率来计算性能比较，那么计算结果将更加准确。

\item \textbf{检查流动性机制。}

将投资者可获得的流动性与基础资产的流动性进行对比。找出在何种情况下，赎回操作、市场影响或ETF的溢价与折价可能会变得重要。

\item \textbf{请检查推荐逻辑是否正确。}

确定该推荐是基于投资者特征还是基于未提供的未来结果。具体来说，需要明确关于投资期限、风险承受能力、流动性需求、税务状况或预期回报的假设。

\end{enumerate}

最终报告应当明确区分可验证的基金特征与预测、假设以及需要更多信息的投资者个人判断。在报告中应保留所有LaTeX命令、公式表达、交叉引用以及数字内容。

%========================================================
% Discussion Questions
%========================================================

\TailModule{questions}{讨论题}

\begin{enumerate}

\item 为什么即使投资者和基金经理都更倾向于较高的投资组合回报，委托资产管理仍然会产生本金-代理人问题呢？

\item ETF的创设-赎回机制与开放式共同基金的赎回机制有何不同？为什么这种差异对投资组合的交易具有重要影响？

\item 为什么一个拥有长期负债的养老基金会持有与那些投资者可以每天赎回的共同基金不同的投资组合呢？

\item 即使被动资产管理者不依赖自主的证券选择，他们如何能够影响公司治理？

\item 哪些资产管理功能可以通过自动化来更高效地执行？而在投资管理的哪些方面，仍然需要特别注意激励机制、市场流动性、治理结构以及人类判断的作用呢？

\end{enumerate}

%========================================================
% Applied Problems
%========================================================

\TailModule{appprobs}{应用题}

\begin{enumerate}

\item \textbf{费用与委托执行相关事项。}

该基金年初拥有约2亿元人民币的资金，并实现了9\%的毛回报。管理费为初始资产的1.2\%。绩效费相当于毛回报超过6\%部分的15\%，该部分是基于初始资产计算的，且在支付管理费之前进行计算的。忽略其他现金流和成本。

\begin{enumerate}

\item 计算管理费用。

\item 计算绩效费用。

\item 计算扣除各项费用后的最终资产价值。

\item 计算投资者简化的净回报。

\end{enumerate}

\item \textbf{共同基金资产净值与赎回事宜。}

该开放式基金的资产为人民币8.4亿元，负债为人民币1200万元，存续股份数量为7200万股。

\begin{enumerate}

\item 计算每股的净资产值。

\item 一位投资者以每股净资产的价格赎回了60万股股票。计算所需的现金数额，同时忽略相关费用和市场影响。

\item 假设该基金为了满足赎回需求而产生了约60万元人民币的交易成本，这些成本由剩余的资金承担。请计算出最终的每股净资产值，并说明谁应承担这些成本。

\item 请选出一种能够改变此类赎回成本分配方式的基金设计工具。

\end{enumerate}

\item \textbf{ETF溢价与折价。}

该ETF的每股市值约为4.00人民币，实际交易价格为4.08人民币。

\begin{enumerate}

\item 计算相对于净资产值的百分比溢价。

\item 描述在溢价足够覆盖交易和融资成本的情况下，潜在的创建交易的过程。

\item 解释为什么在基础证券流动性不足或市场关闭的情况下，溢价或折扣并不需要立即消失。

\end{enumerate}

\item \textbf{养老金资金筹措与负债敏感性。}

养老基金的资产为人民币950亿元，预计负债的现值为人民币1000亿元。

\begin{enumerate}

\item 计算资金比例。

\item 如果负债额上升至1060亿元人民币，而资产保持不变，那么需要计算新的资金比率。

\item 解释长期贴现率下降如何增加长期负债的现值。

\end{enumerate}

\item \textbf{机构交易与实施缺口。}

有一家基金会决定在参考价格为20元人民币时购买500,000股股票。平均执行价格约为20.12元人民币。

\begin{enumerate}

\item 使用LaTeX命令、公式、交叉引用以及数字来表示计算结果。  
根据方程~\ref{eq:ch14-implementation-shortfall}来计算实施过程中的不足，同时忽略相关费用的影响。

\item 将短缺金额表示为订单参考价值的百分比。

\item 解释快速完成订单与降低市场影响之间的权衡关系。

\end{enumerate}

\end{enumerate}

%========================================================
% Data Application
%========================================================

\TailModule{dataapp}{数据应用}

\textbf{在中国，公开募集的资金、ETF以及被动投资策略}

使用中国证券投资基金业协会、上海证券交易所、深圳证券交易所、基金招募说明书和定期基金报告等官方或原始来源的数据。选择少量可在共同期间比较的主动管理基金，以及指数基金或~ETF。

首先，需要提供有关基金规模、认购或赎回情况的信息，以及管理费用的相关数据。对于ETF来说，还需要收集市场价格以及选定日期的净资产值或预估净资产值，然后利用公式~\ref{eq:ch14-etf-premium-discount}来计算溢价或折扣。同时，要明确说明市场价格与净资产值指标在时间上的对应关系。

其次，将各被动基金的回报与它们所设定的基准指标进行长期比较，并计算其跟踪差异。利用基金披露信息来分析费用、现金持有量、抽样方法、重新平衡操作等因素是否有助于解释这些差异。不过，不要将过去的跟踪表现视为未来表现的预测指标。

第三，检查所选基金如何行使所有权。在信息披露可用的前提下，可以使用投票记录、定期报告或其他与管理相关的信息来评估被动投资组合的构建是否伴随着积极的投票或参与行为。

请列出各基金的名称、份额类别、样本周期、基准定义、费用规定、数据来源，以及针对收益或净资产值的任何调整措施。务必区分实际的数据与关于投资者需求、公司治理或市场影响的推测性描述。

%========================================================
% Further Reading
%========================================================

\TailModule{furtherread}{延伸阅读}

\begin{itemize}

\item Chevalier, Judith, and Glenn Ellison (1997),

``Risk Taking by Mutual Funds as a Response to Incentives（共同基金对激励的风险承担反应）,''

\emph{Journal of Political Economy}, 105(6), 1167--1200.

\item Sirri, Erik R., and Peter Tufano (1998),

``Costly Search and Mutual Fund Flows（有成本的搜寻与共同基金资金流动）,''

\emph{The Journal of Finance}, 53(5), 1589--1622.

\item Berk, Jonathan B., and Richard C. Green (2004),

``Mutual Fund Flows and Performance in Rational Markets（理性市场中的共同基金资金流动与业绩）,''

\emph{Journal of Political Economy}, 112(6), 1269--1295.

\item Appel, Ian R., Todd A. Gormley, and Donald B. Keim (2016),

``Passive Investors, Not Passive Owners（被动投资者，并非被动所有者）,''

\emph{Journal of Financial Economics}, 121(1), 111--141.

\item Ben-David, Itzhak, Francesco Franzoni, and Rabih Moussawi (2018),

``Do ETFs Increase Volatility?（ETF是否增加波动？）''

\emph{The Journal of Finance}, 73(6), 2471--2535.

\end{itemize}

\ChapterTail{summary,keyterms,aiprompt,questions,appprobs,dataapp,furtherread}
